% Generated draft: requires mathematical and source review.
\documentclass[11pt]{article}
\usepackage[margin=25mm]{geometry}
\usepackage{fontspec}
\setmainfont{Latin Modern Roman}
\usepackage{amsmath,amssymb,mathtools,unicode-math}
\setmathfont{Latin Modern Math}
\usepackage{enumitem,needspace,xurl,hyperref}
\setlength{\parindent}{0pt}
\setlength{\parskip}{5pt}
\setlength{\emergencystretch}{3em}
\newcommand{\N}{\mathbb{N}}
\newcommand{\Z}{\mathbb{Z}}
\newcommand{\Q}{\mathbb{Q}}
\newcommand{\R}{\mathbb{R}}
\newcommand{\C}{\mathbb{C}}
\newcommand{\F}{\mathbb{F}}
\newcommand{\A}{\mathbb{A}}
\newcommand{\PP}{\mathbb{P}}
\newcommand{\E}{\mathbb{E}}
\newcommand{\eps}{\varepsilon}
\newcommand{\dd}{\,\mathrm{d}}
\begin{document}
\section*{UnsolvedMath ID: 700932}
\textbf{Source problem code:} catalogue entry
\subsection{Definitions and mathematical statement}
Work in the language of groups, with multiplication, inversion, and identity element \(1\). Write \([x,y]=x^{-1}y^{-1}xy\). Let \(\mathcal K\) be the class of all groups \(G\) satisfying
\[
\forall x,y,z\in G,\qquad [[x,y],z]=1,
\]
and, for every integer \(n\geq 1\),
\[
\forall x\in G,\qquad x^n=1\Longrightarrow x=1.
\]
Thus \(\mathcal K\) consists of the torsion-free nilpotent groups of nilpotency class at most \(2\), including abelian groups and the trivial group. No finiteness or finite-generation hypothesis is imposed.

A group quasi-identity is a universally quantified implication
\[
\forall x_1,\ldots,x_r,\qquad
\left(\bigwedge_{i=1}^{m}s_i(x_1,\ldots,x_r)=t_i(x_1,\ldots,x_r)\right)
\Longrightarrow s(x_1,\ldots,x_r)=t(x_1,\ldots,x_r),
\]
where all displayed expressions are group terms and \(m\) is finite and may be zero. A quasivariety of groups is the class of all groups satisfying some set of group quasi-identities. The phrase ``quasivarieties of nilpotent torsion-free groups of nilpotency class at most \(2\)'' is interpreted as quasivarieties of groups contained in \(\mathcal K\).

Let \(\mathcal L\) be the lattice of these quasivarieties, ordered by inclusion. Since \(\mathcal K\) is itself axiomatized by quasi-identities, its lattice operations are
\[
\mathcal U\wedge\mathcal V=\mathcal U\cap\mathcal V,
\qquad
\mathcal U\vee\mathcal V=
\bigcap\{\mathcal W\in\mathcal L:\mathcal U\cup\mathcal V\subseteq\mathcal W\}.
\]
The join is therefore the smallest quasivariety containing both classes.

\textbf{Question (A. I. Budkin).} Is \(\mathcal L\) distributive? Equivalently, does
\[
\mathcal U\cap(\mathcal V\vee\mathcal W)
=(\mathcal U\cap\mathcal V)\vee(\mathcal U\cap\mathcal W)
\]
hold for every \(\mathcal U,\mathcal V,\mathcal W\in\mathcal L\)?
\subsection{Short English statement}
Consider all quasivarieties contained in the class of torsion-free groups of nilpotency class at most \(2\). Ordered by inclusion, they form a lattice whose meet is intersection and whose join is the quasivariety generated by the two classes. Budkin asks whether this lattice is distributive.
\subsection{Sources}
\begin{itemize}
\item Local UnsolvedMath registry, permanent ID \(700932\), exact source problem code \texttt{catalogue entry}, title ``Kourovka Notebook Problem 13.14.'' The supplied record labels the problem open; this label is supplied source metadata, not an independently verified mathematical or literature assessment.
\item \emph{The Kourovka Notebook (Unsolved Problems in Group Theory)}, original problem \(13.14\), proposed by A. I. Budkin in the 13th issue (1995). The supplied background also identifies the 21st issue edition (2026), \texttt{arXiv:1401.0300v48}. The supplied collection description names E. I. Khukhro and V. D. Mazurov as editors. These bibliographic details are reproduced from the record and have not been independently verified.
\end{itemize}
\end{document}
